%%%%%%%%%%%%%%%%%%%%%%%%%%%%%%%%%%%%%%%%%%%%%%%%%%%%%%%%%%%%%%%%%%%%%%%%%%%%%%
%%%%%%%%% Ajit Kumar Sahoo, SCIS, University of  Hyderabad,
%%%%%%%%%  (email:sahooajitkumar85@gmail.com)
%%%%%%%%% https://ajitsahoocs.github.io/
%%%%%%%%%%%%%%%%%%%%%%%%%%%%%%%%%%%%%%%%%%%%%%%%%%%%%%%%%%%%%%%%%%%%%%%%%%%%%%%%%%%%%%%%%%%%%%
% Beamer Presentation
% LaTeX Template
\documentclass{beamer}
\hypersetup{pdfpagemode=FullScreen} %full screen mode option 
\setbeamertemplate{navigation symbols}{} % To remove navigation panel on the lower right corner
\usepackage[english]{babel} %designed for typesetting lebels(addr lebel,sticky lebels,etc)
\usepackage[latin1]{inputenc} %Accept different input encodings
\usepackage{amsmath,amssymb} 
\usepackage[T1]{fontenc} %for font encoding
\usepackage{times} % use times font instead of default
\usepackage{curves} %for drawing curves
\usepackage{verbatim} %paragraph making environment 
\usepackage{multimedia} %for multimedia like animation,movie etc...
\usepackage{mathptmx} % font style
\usepackage{graphicx} % Allows including images
\usepackage{booktabs} % Allows the use of \toprule, \midrule and \bottomrule in tables
\usepackage{hyperref}
\usepackage{xcolor}
\usepackage{subfig}

\usepackage{algorithm,algorithmic}
\renewcommand{\algorithmicrequire}{\textbf{Input:}}
\renewcommand{\algorithmicensure}{\textbf{Output:}}
\usepackage{lipsum}
\setbeamertemplate{caption}[numbered] % For numbering figures


\setbeamertemplate{frametitle continuation}{}
\mode<presentation> {

% The Beamer class comes with a number of default slide themes
% which change the colors and layouts of slides. Below this is a list
% of all the themes, uncomment each in turn to see what they look like.

%\usetheme{default} %default
\usetheme{Madrid} % popular one
%\usetheme[secheader]{Madrid}
%\usetheme{CambridgeUS}
%%%%%%%%%%%%%%%%%%%%%%%%%%%%%%%
%\usetheme{AnnArbor}
%\usetheme{Antibes}
%\usetheme{Bergen}
%\usetheme{Berkeley}
%\usetheme{Berlin}
%\usetheme{Boadilla}
%\usetheme{CambridgeUS}
%\usetheme{Copenhagen}
%\usetheme{Darmstadt}
%\usetheme{Dresden}
%\usetheme{Frankfurt}
%\usetheme{Goettingen}
%\usetheme{Hannover}
%\usetheme{Ilmenau}
%\usetheme{JuanLesPins}
%\usetheme{Luebeck}
%\usetheme{Malmoe}
%\usetheme{Marburg}
%\usetheme{Montpellier}
%\usetheme{PaloAlto}
%\usetheme{Pittsburgh}
%\usetheme{Rochester}
%\usetheme{Singapore}
%\usetheme{Szeged}
%\usetheme{Warsaw}
%%%%%%%%%%%%%%%%%%%%%%%%%%%%%%%%%%%%%%%%%%%%%%%%%%%%%%%%%%%%%%%%%%%%%%%%%%%%%%%%
% As well as themes, the Beamer class has a number of color themes
% for any slide theme. Uncomment each of these in turn to see how it
% changes the colors of your current slide theme.
%\usecolortheme{dolphin} %default
%\usecolortheme{dolphin}
%\usecolortheme{orchid} % best one blue
%\usecolortheme{dove}
%\usecolortheme{crane} %yellow color
%\usecolortheme{beaver} %grey
%\usecolortheme{whale} %blue


%\usecolortheme{albatross}
%\usecolortheme{beetle}
%\usecolortheme{fly}
%\usecolortheme{lily}
\usecolortheme{rose}
%\usecolortheme{seagull}
%\usecolortheme{seahorse}
%\usecolortheme{wolverine}

%\setbeamertemplate{footline} % To remove the footer line in all slides uncomment this line
%\setbeamertemplate{footline}[page number] % To replace the footer line in all slides with a simple slide count uncomment this line

%\setbeamertemplate{navigation symbols}{} % To remove the navigation symbols from the bottom of all slides uncomment this line
}

\usepackage{graphicx} % Allows including images
\usepackage{booktabs} % Allows the use of \toprule, \midrule and \bottomrule in tables
\setbeamercovered{transparent}
\setbeamertemplate{bibliography item}[text]
\setbeamertemplate{theorems}[numbered]
\setbeamerfont{title}{size=\Large}%\miniscule,command,tiny, scriptsize,footnotesize,small,normalsize,large,Large,LARGE,huge,Huge,HUGE
\setbeamerfont{date}{size=\tiny }%{\fontsize{40}{48} \selectfont Text}

\setbeamertemplate{itemize items}[ball] % if you want a ball
\setbeamertemplate{itemize subitem}[circle] % if you wnat a circle
\setbeamertemplate{itemize subsubitem}[triangle] % if you want a triangle
%%%%%%%%%%%%%%%%%%%%%%%%%%%%%%%%%%%%%%%%%%%%%%%%%%%%%%%%%%%%%%%%%%%%%%%%%%%%%%%%%%%%%%%%%%
%----------------------------------------------------------------------------------------
%	TITLE PAGE
%----------------------------------------------------------------------------------------
%%%%%%%%%%%%%%%%%%%%%%%%%%%%%%%%%%%%%%%%%%%%%%%%%%%%%%%%%%%%%%%%%%%%%%%%%%%%%%%%%%%%%%%%%%
\title[\textit{Ph.D. Viva Voce Presentation} ]{Title of the thesis or dissertation, Title of the thesis or dissertation  }
%\title[ML-intelligent sensing]{Non-contact intelligent ultrasonic  sensing  using machine learning approaches} % The short title appears at the bottom of every slide, the full title is only on the title page

%\author[AKS]{Ajit Kumar Sahoo \texorpdfstring{\scriptsize Regd No: 16MCPC03}{}}
\author[AKS]{\texorpdfstring {\scriptsize \textit{Ph.D.  Thesis Defence Presentation \\ \vspace{0.1cm} }} \texorpdfstring{\scriptsize by \\ \vspace{0.1cm}} \texorpdfstring{\normalsize Name {}}{}}

\institute[UoH] % Your institution as it will appear on the bottom of every slide, may be shorthand to save space
{

{\small \textit{Supervisor : }}
\medskip
{\small Name}
\vspace{0.2 cm}
\texorpdfstring{\\ \scriptsize \textit{DAC Committee  : }}
\medskip
{\small Names\\}
%\begin{center}
%\includegraphics[width=0.1\textwidth]{uoh.png}
%\end{center}
\vspace{1.8cm}
\medskip
School of Computer and Information Sciences\\
University of Hyderabad, Hyderabad, India, 500046% Your institution for the title page

%\medskip
%\textit{john@smith.com} % Your email address
}
%This will place the image at position "30 right/left and 120 up/down" relative to the top left corner of the current page.
\titlegraphic{%
  \begin{picture}(0,0)
    \put(30,124){\makebox(0,0)[rt]{\includegraphics[width=2cm]{uoh.png}}}
   % \put(20,105){\makebox[10]{\includegraphics[width=1.3cm]{uoh.png}}}
  \end{picture}}
% \vspace*{-\baselineskip} 
 %\vspace*{-\baselineskip} 
 
 \date{\today} % Date, can be changed to a custom date
%\date{}

%%%%%%%%%%%%%%%%%%%%%%%%%%%%%%%%%%%%%%%%%%%%%%%%%%%%%%%%%%%%%%%%%%%%%%%%%%%%%%%%%%%%%%%%%%%%%%%%%%%%%%%%%%%%%%%

\begin{document}

\begin{frame}
\titlepage % Print the title page as the first slide\\
\end{frame}
\begin{frame}
  \frametitle{Contents}
 % \tableofcontents[pausesections,shaded]
 \tableofcontents
\end{frame}
% TABLE OF CONTENTS AT BEGIN OF EACH SECTION
\AtBeginSection[]{
  \begin{frame}<beamer>
    \frametitle{Current Section}
   \tableofcontents[currentsection]
  \end{frame}}
%----------------------------------------------------------------------------------------
%	PRESENTATION SLIDES
%----------------------------------------------------------------------------------------
\section{Introduction}
\begin{frame}[allowframebreaks]
\frametitle{Introduction}

   \begin{itemize}
       
       \item statement 1
       \item statement 2
       \item statement 3
   \end{itemize}


  \begin{figure} %htbp! %https://www1.cmc.edu/pages/faculty/aaksoy/latex/latexfive.html
    %\includegraphics[width=50mm,scale=0.5]{}
    %\includegraphics[trim=left bottom right top, clip]{file}
    %\includegraphics[trim=1200 400 1500 1500,clip,width=\textwidth]{crab}
	\includegraphics[trim=2cm 2cm 2cm 1cm,clip,width=\textwidth]{images/intelligent_sensor_diagram.pdf}
	%\includegraphics[scale=0.5]{figures/ch1/intelligent_sensor.pdf}
	\caption{Block diagram of intelligent sensor }
	\label{figure1}
\end{figure}

\end{frame}

\subsection{Motivation}
\begin{frame}{Motivation}
     \begin{itemize}
       
       \item statement 1
       \item statement 2
       \item statement 3
   \end{itemize}
\end{frame}

\subsection{Research Gaps}
\begin{frame}{Research Gaps}
     \begin{itemize}
       
       \item statement 1
       \item statement 2
       \item statement 3
   \end{itemize}
\end{frame}

\subsection{Assumptions}
\begin{frame}{Assumptions}
     \begin{itemize}
       
       \item statement 1
       \item statement 2
       \item statement 3
   \end{itemize}
\end{frame}





\begin{frame}{Characteristics of ultrasonic wave}
    
         \begin{equation*}
      c = \sqrt{\frac{\gamma P}{\rho}} 
         \end{equation*}
       \begin{equation}
         = \sqrt{\frac{\gamma RT}{M}}
      \end{equation}
         
where $\gamma$ is the adiabatic index (ratio of specific heats), $R$ is the universal gas constant, $T$ is the temperature in kelvin, and $M$ is the molecular mass. P is the pressure\footnotemark. 
 \footnotetext[1]{\textit{A. K. Sahoo and S. K. Udgata, "A Novel ANN-Based Adaptive Ultrasonic Measurement System for Accurate Water Level Monitoring," in IEEE Transactions on Instrumentation and Measurement, vol. 69, no. 6, pp. 3359-3369, June 2020.}}      
\end{frame}
%------------------------------------------------------------------------------------------------------------------------------------------
\section{Background and Literature Survey} 
\begin{frame}[allowframebreaks]
   \frametitle{Background and Literature Survey}


  \begin{itemize}
       
       \item statement 1
       \item statement 2
       \item statement 3
   \end{itemize}


\end{frame}

%------------------------------------------------------------------------------------------------------------------------------------------
\section{Research Objectives} 
\begin{frame}[allowframebreaks]
   \frametitle{Research Objectives}


  \begin{itemize}
       
       \item statement 1
       \item statement 2
       \item statement 3
   \end{itemize}


\end{frame}


\begin{frame}[allowframebreaks]
\frametitle{Main Contributions}
\begin{itemize}
   \item statement 1
    \item statement 2
    \item statement 3
\end{itemize}
\end{frame}







%-------------------------------------------------------------------------------------------------------------------------------
\section{Contribution-1: description }
\begin{frame} [allowframebreaks]
\frametitle{Analog Echo Signal }

\begin{definition} 
        This is a definition.
    \end{definition}
    
    \begin{theorem} 
        This is a theorem. 
    \end{theorem}
    
    \begin{lemma} 
        This is a proof idea.
    \end{lemma}


\end{frame}






\begin{frame}{Methodology}

\begin{algorithm}[H]
\begin{algorithmic}[1]
\REQUIRE input data
\ENSURE output data 
\FOR{$j=1$ to $N$}
\STATE Statement
\ENDFOR
\end{algorithmic}
\caption{pseudocode for algorithms}
\label{alg1}
\end{algorithm}


\end{frame}






    

\begin{frame}{Experiment}
    \begin{block}{Standard Block}
        This is a standard block.
    \end{block}
    \begin{alertblock}{Alert Message}
        This block presents alert message.
    \end{alertblock}
    \begin{exampleblock}{An example of typesetting tool}
        Example: MS Word, \LaTeX{}
    \end{exampleblock}
  
\end{frame}



\begin{frame}{Result}
    \begin{proof} 
        This is a proof. 
    \end{proof}
    
    \begin{corollary}
        This is a corollary
    \end{corollary}
    
    \begin{example}
        This is an example 
    \end{example}
    
\end{frame}

\begin{frame}{Conclusions}
\begin{itemize}
   \item statement 1
    \item statement 2
    \item statement 3
\end{itemize}
\end{frame}

%-------------------------------------------------------------------------------------------------------------------------------
\section{Contribution-2: description }
\begin{frame} [allowframebreaks]
\frametitle{Contribution-2}

\begin{itemize}
    \item statement 1
    \item statement 2
    \item statement 3
\end{itemize}


\end{frame}



%-------------------------------------------------------------------------------------------------------------------------------
\section{Contribution-3: description }
\begin{frame} [allowframebreaks]
\frametitle{Contribution-3}

\begin{itemize}
    \item statement 1
    \item statement 2
    \item statement 3
\end{itemize}


\end{frame}



%-------------------------------------------------------------------------------------------------------------------------------
\section{Contribution-4: description }
\begin{frame} [allowframebreaks]
\frametitle{Contribution-4}

\begin{itemize}
    \item statement 1
    \item statement 2
    \item statement 3
\end{itemize}


\end{frame}


%-------------------------------------------------------------------------------------------------------------------------------

\section{Conclusions and Future Scope}
\begin{frame}[allowframebreaks]
\frametitle{Conclusions}
\begin{itemize}
    \item statement 1
    \item statement 2
    \item statement 3
\end{itemize}
\end{frame}



\begin{frame}[allowframebreaks]
\frametitle{Future Scope}
While concluding the thesis, we also identified a few important scopes for future work and intended to carry forward the work in the following areas;
\begin{itemize}
    \item statement 1
    \item statement 2
    \item statement 3
\end{itemize}
 
\end{frame}
%-------------------------------------------------------------------------------------------------------------------------------
\section{Reviewers Comments and Response }
\begin{frame}[allowframebreaks]
\frametitle{Reviewers Comments and Reply}
\begin{enumerate}
\item Question 1 
   \begin{itemize}
       \item Reply 1:
   \end{itemize}
\item Question 2
      \begin{itemize}
       \item Reply 2:
   \end{itemize}
\item Question 3
     \begin{itemize}
       \item Reply 3:
   \end{itemize}
\item Question 4
     \begin{itemize}
       \item Reply 4:
   \end{itemize}
\end{enumerate}

\end{frame}




%-------------------------------------------------------------------------------------------------------------------------------
\section{References}
\begin{frame}[allowframebreaks]
\frametitle{Important References}
 %\nocite{david2017fundamentals}
 \nocite{*}
 \bibliographystyle{ieeetr}
\bibliography{References}
\end{frame}



%--------------------------------------------------------------------------------------------------------------------------------
\section{Publications}
\begin{frame} [allowframebreaks]
\frametitle{Publications}
\begin{enumerate}
\item publication 1 
\item publication 2
\item publication 3
\item publication 4
\end{enumerate}

\end{frame}


%--------------------------------------------------------------------------------------------------------------------------------
\section{Acknowledgement}
\begin{frame} [allowframebreaks]
\frametitle{Acknowledgement}
\lipsum[1]


\end{frame}


%-------------------------------------------------------------------------------------------------------------------------------

%-------------------------------------------------------------------------------------------------------------------------------

%\section{Information processing  issues}
%-----------------------------------------------
%\begin{frame}
%\frametitle{Information processing issues}
%\begin{itemize} 
%\item In collaborative processing
%\begin{itemize}
%\item Target detection, localization, communication, storage, querying...
%\end{itemize}
%\color{lightgray}
%\item  In networking
%\begin{itemize}
%\color{lightgray}
%\item  aggregation, routing, Data naming,...
%\end{itemize}
%\end{itemize}
%\end{frame}
%-----------------------------------------------
%\begin{frame}[fragile] % Need to use the fragile option when verbatim is used in the slide
%\frametitle{Citation}
%An example of the \verb|\cite| command to cite within the presentation:\\~
%\end{frame}
%------------------------------------------------


%-----------------------------------------------------------------------------------------

\begin{frame}
\Huge{\centerline{Thank You}}
\end{frame}

%----------------------------------------------------------------------------------------

\end{document}